\chapter{Toán tử và cơ chế suy luận}\label{chap:reasoning}

\section{Chuẩn hóa và trích đặc trưng}
Gọi $N(\ell)$ là phép chuẩn hóa nhãn được cài đặt trong
\src{normalize_label}. Hàm thực hiện lần lượt Unicode NFKC, chuyển dạng chữ
bằng \texttt{casefold}, thay một tập dấu phân cách thành khoảng trắng,
rút gọn khoảng trắng và thay từ theo bảng bí danh.
Bảng bí danh hiện chỉ có ba mục: \texttt{haemorrhage} thành
\texttt{hemorrhage}, \texttt{haemorrhagic} thành \texttt{hemorrhagic}
và \texttt{renal} thành \texttt{kidney}.
Đây là tri thức từ vựng được khai báo thủ công, không phải mô hình học đồng nghĩa.

Tập từ nội dung được định nghĩa theo code là
\begin{equation}
  T(\ell)=\operatorname{set}(\operatorname{split}(N(\ell)))\setminus W,
  \quad
  W=\{\text{a, an, and, due, of, or, the, to}\}.
  \label{eq:tokens}
\end{equation}
Biểu diễn tập hợp bỏ qua thứ tự và số lần xuất hiện của từ.
Phép chuẩn hóa mã là loại khoảng trắng rồi viết hoa; dấu chấm trong mã
được giữ. Tuy nhiên, chuẩn hóa mã chủ yếu phục vụ tìm kiếm và kiểm tra
self-mapping, không được áp dụng thống nhất ở mọi bước tạo định danh.

\section{Phát hiện xung đột theo cụm từ}
Gọi $D$ là tập các cặp cụm từ đối lập được khai báo.
Một phần tử điển hình là $(\text{acute},\text{chronic})$;
những cặp khác gồm benign và malignant; left và right; type 1 và type 2;
with hemorrhage và without hemorrhage; hemorrhagic và without hemorrhage;
with perforation và without perforation.
Phép kiểm tra có dạng
\begin{equation}
\begin{split}
X(\ell_s,\ell_t)=\{(u,v)\in D:\;&
[u\sqsubseteq N(\ell_s)\land v\sqsubseteq N(\ell_t)]\\
&{}\lor[v\sqsubseteq N(\ell_s)\land u\sqsubseteq N(\ell_t)]\},
\end{split}
\label{eq:conflict}
\end{equation}
trong đó $\sqsubseteq$ là quan hệ xuất hiện như \emph{chuỗi con}.
Ký hiệu này không biểu diễn quan hệ chuyên biệt hóa bệnh.
Điều kiện xung đột là $X(\ell_s,\ell_t)\ne\varnothing$.

Do dùng chuỗi con, phép kiểm tra chưa hiểu phạm vi phủ định hoặc cấu trúc
“unspecified as acute or chronic”. Danh sách xung đột được lưu vào proof
giúp người dùng phát hiện những nhận định quá đơn giản này.

\section{Năm luật và thứ tự áp dụng}
Đặt $S=T(\ell_s)$, $T=T(\ell_t)$ và
$U=\{\text{unspecified},\text{other}\}$.
Các điều kiện $P_j$ của năm luật được viết tương ứng với Python như sau:
\begin{align}
P_1(m)&:\ N(\ell_s)=N(\ell_t), & y_1&=A, \label{eq:r1}\\
P_2(m)&:\ X(\ell_s,\ell_t)\ne\varnothing, & y_2&=C, \label{eq:r2}\\
P_3(m)&:\ X=\varnothing \land S\ne\varnothing\land T\ne\varnothing
            \land(S\subsetneq T\lor T\subsetneq S), & y_3&=B, \label{eq:r3}\\
P_4(m)&:\ X=\varnothing \land S\cap T\ne\varnothing
            \land(S\cup T)\cap U\ne\varnothing, & y_4&=B, \label{eq:r4}\\
P_5(m)&:\ \mathrm{true}, & y_5&=\bot. \label{eq:r5}
\end{align}
Trong \eqref{eq:r3}--\eqref{eq:r4}, $X$ viết gọn cho
$X(\ell_s,\ell_t)$. Luật thứ tư chỉ yêu cầu có từ chung và có từ chỉ tính
khái quát; nó không kiểm tra rằng từ chung là tên bệnh hay rằng hai
khái niệm thuộc cùng một nhánh ontology.

\begin{table}[htbp]
\centering\small
\caption{Mã luật và bằng chứng được trả về khi luật khớp}\label{tab:rules}
\begin{tabularx}{\textwidth}{c p{6.0cm} Y}
\toprule
Ưu tiên & Mã luật & Premise trong proof\\
\midrule
1 & \src{R-EXACT-LABEL} & Hai nhãn sau chuẩn hóa.\\
2 & \src{R-EXPLICIT-QUALIFIER-CONFLICT} & Hai nhãn gốc và danh sách xung đột.\\
3 & \src{R-COMPATIBLE-SPECIFICITY} & Hai nhãn gốc và explicitConflict=false.\\
4 & \src{R-GENERALIZED-CONDITION} & Hai nhãn gốc và explicitConflict=false.\\
5 & \src{R-INSUFFICIENT-EVIDENCE} & Hai nhãn gốc.\\
\bottomrule
\end{tabularx}
\end{table}

\src{COKBReasoner.assess} dừng ở luật đầu tiên khớp:
\begin{equation}
 j^*(m)=\min\{j:P_j(m)=\mathrm{true}\},\qquad
 f_{\K}(m)=(y_{j^*},\texttt{rule\_based},\pi_{j^*}).
 \label{eq:firstmatch}
\end{equation}
Luật mặc định $P_5$ bảo đảm luôn có đầu ra với tập luật chuẩn.
Thứ tự là một phần ngữ nghĩa của hệ thống: nhãn giống hệt nhau được
xét trước xung đột. Một tập luật tùy chỉnh không có fallback có thể làm
reasoner phát sinh lỗi khi không luật nào khớp.

\noindent\begin{minipage}{\linewidth}
\begin{lstlisting}[language=Python,caption={Giả mã tương ứng với vòng lặp của COKBReasoner.assess},label={lst:reason}]
def assess(mapping, ordered_rules):
    for rule in ordered_rules:
        match = rule.evaluate(mapping)
        if match is None:
            continue
        proof = build_proof(mapping, match)
        return build_assessment(mapping, match, proof)
    raise RuntimeError("Rule set did not produce an assessment")
\end{lstlisting}
\end{minipage}

Giả mã trong Mã minh họa~\ref{lst:reason} lược bỏ thao tác dựng dataclass.
Đây là cơ chế áp dụng luật một lượt trên từng mapping, không lặp đến điểm
cố định, không sử dụng assessment của mapping trước làm tiền đề cho mapping
sau và không tự suy diễn bắc cầu giữa các hệ mã.

\section{Vết suy luận và ý nghĩa của tính tái lập}
Với một luật khớp, vết suy luận có cấu trúc
\begin{equation}
  \pi=(q,\mathcal{P},r,\mathcal{E}),
  \qquad
  \frac{\mathcal{P}\ ;\ r}{q},
  \label{eq:proof}
\end{equation}
trong đó $q$ là kết luận, $\mathcal{P}$ là tập premise, $r$ là mã luật và
$\mathcal{E}$ là nguồn bằng chứng. Ký hiệu phân số trong
\eqref{eq:proof} mô tả cấu trúc giải trình, không khẳng định chương trình
đã có một bộ kiểm chứng chứng minh độc lập.

Proof của luật 3 và 4 lưu nhãn nhưng chưa lưu tường minh tập từ sau xử lý,
quan hệ bao hàm hay từ chung. Proof của fallback cũng chưa liệt kê vì sao
từng luật trước thất bại. Do đó, vết hiện tại đủ để xác định luật và dữ
liệu đầu vào, nhưng chưa là nhật ký đầy đủ của mọi bước tính toán.

Tính xác định được hiểu là cùng input, cùng code và cùng thứ tự luật thì
thu được cùng assessment. Trong \src{COKBRepository.load}, nếu JSON đã
có assessments, code tính lại chúng bằng luật hiện hành.
Vì thế, mở một store sau khi sửa luật có thể cho kết quả mới; muốn tái lập
lịch sử cần lưu phiên bản code bên cạnh dữ liệu.

\section{Độ phức tạp và giới hạn biểu diễn}
Gọi $n$ là số mapping, $k$ là số luật và $L$ là tổng độ dài hai nhãn.
Với số cặp xung đột cố định, mỗi phép chuẩn hóa và kiểm tra có chi phí
xấp xỉ tuyến tính theo $L$. Tổng chi phí đánh giá có dạng
\begin{equation}
  O(nkL),\qquad k=5 \text{ trong tập luật mặc định}.
\end{equation}
Một số phép chuẩn hóa được lặp lại qua các luật; chưa có cache đặc trưng.
Truy vấn \src{find_mappings} quét danh sách mapping nên có chi phí
$O(n)$ nếu xem chiều dài mã là bị chặn.
Các công thức này là phân tích thuật toán từ source, không phải kết quả
đo tốc độ hay cam kết hiệu năng trên graph đầy đủ.
